\textbf{Scale.} This is a small CPU study (about 30 minutes of wall-clock in total), with one modulus ($p=23$), one architecture and 3--5 seeds per cell; all conclusions are at that scale.

\textbf{Tasks and splits.} Modular addition mod 23; a random train/held-out split of the 529 pairs per seed, training fraction $f\in\{0.4,0.5,0.7\}$; the held-out set is every pair not in the training set. The \emph{main cell} is $\lambda=1$, $f=0.5$, with seeds $0$--$4$ for the three systems. The \emph{grid} crosses $\lambda\in\{0,1,3\}$ with $f\in\{0.4,0.5,0.7\}$ (the main cell is not repeated), seeds $0$--$2$; in the run registry these rows were relabelled after the runs into group \texttt{main} (see \texttt{experiments/PROTOCOL.md}). Two extra seeds (5, 6) of the main cell, kept for the curves, are reported in Table~\ref{tab:pooled}. Two sensitivity groups run at the main cell with seeds $0$--$4$: curriculum length $T_c\in\{250,500,2000,4000\}$ (the default $T_c=1000$ is in the main group) and the shuffled-order control.

\textbf{Hyper-parameters and tuning budget.} Model and optimiser as in Section~\ref{sec:method}; budget 4000 steps; evaluation every 25 steps. The learning rate and the main cell were chosen from about ten pilot runs of uniform sampling with a separate seed (100) that appears in no table and is not in the run registry (the pilot was exploratory and its outputs were not kept, so it cannot be re-checked); it led us to $\eta=3\times10^{-3}$, $\lambda=1$, $f=0.5$. No hyper-parameter was tuned per system, which favours uniform sampling (the recipe the pilot used) over the orderings.

\textbf{Statistics.} Runs that never reach 95\% are censored at 4000 steps, so mean and standard deviation of steps\_to\_95 are lower bounds; we additionally report the fraction of seeds that reached 95\%. Significance tests are those of the harness comparison tool (Welch and paired tests of per-seed values against the curriculum row), with only 5 (or 3) seeds per cell, so they have little power.

\textbf{Hardware and cost.} Shared Apple-silicon CPU, 2 threads per process, at most two processes in parallel, PyTorch; a run takes at most about 30\,s. Code, logs and the run registry are in the repository linked in the author block.
